\subsection{Domain of a Rational Function} 
Recall that we get a rational function by dividing two polynomial functions, and that the domain of a polynomial function is all real numbers. This tells us how to find the domain of any rational function.
\begin{theorem}[Domain of a Rational Function]
The domain of a rational function $R(x)=\dfrac{P(x)}{Q(x)}$ is $\Set{x\mid Q(x)\neq 0}$.
\end{theorem}
We can also see this directly: we can plug any number we like into a polynomial, but we can't \makeblank[1in] by \makeblanksmall.
\begin{Example}
Find the domain of the rational function $f(x)=\dfrac{3x-5}{2x+1}$.
\begin{lpic}{}{4}
\regtextleft{0}{-1}{To do this we must solve };
\regtextleft{0}{-4}{$\dom f={}$};
\end{lpic}
\end{Example}
\begin{Example}
Find the domain of the rational function $g(x)=\dfrac{2x-3}{x^2-7x+12}$.
\begin{lpic}{}{5}
\regtextleft{0}{-1}{To do this we must solve };
\regtextleft{0}{-5}{$\dom g={}$};
\end{lpic}
\end{Example}
